\documentclass[12pt]{article}
\usepackage[utf8]{inputenc}
\usepackage[T1]{fontenc}
\usepackage[margin=1in]{geometry}
\usepackage{setspace}
\usepackage{xcolor}
\usepackage[colorlinks=true, urlcolor=blue!60!black, linkcolor=blue!60!black]{hyperref}
\setlength{\parindent}{0em}
\setlength{\parskip}{0.45em}
\renewcommand{\baselinestretch}{1.0}
\begin{document}
\begin{center}
{\Large \textit{Shelley v. Kraemer}, 334 U.S. 1 (1948)\par}
\vspace{0.3em}
{\normalsize Case Brief --- One-Page Edition\par}
\end{center}

\textbf{Memory Jogger.} State-court enforcement of racially restrictive covenants is state action denying equal protection (\href{https://supreme.justia.com/cases/federal/us/334/1/}{Justia}).

\textbf{Facts.} Recorded covenants in St. Louis and Detroit barred occupancy by anyone ``not of the Caucasian race.'' The Shelleys and the McGhees, Black families, bought restricted homes from willing sellers; neighbors sued to oust them.

\textbf{Procedural History.} The Missouri and Michigan supreme courts ordered enforcement, divesting the Shelleys of title and ordering the McGhees to vacate.

\textbf{Votes.} Reversed 6--0: Vinson for the Court; Reed, Jackson, and Rutledge did not participate.

\textbf{Holding.} Private covenants alone are lawful, but judicial enforcement against willing buyers and sellers is state action that denies equal protection.

\textbf{Reasoning.} The covenants targeted a class defined ``wholly in terms of race or color'' and would be plainly invalid as statutes under \textit{Buchanan v. Warley}, but the Amendment reaches only state action. Courts act for the State, even when enforcing common-law rules through fair proceedings. But for the courts' ``active intervention,'' the families would have lived in their homes. Equal enforcement against white buyers is no defense, because the rights are ``personal rights.'' \textit{Barrows v. Jackson} later barred damages actions.

\textbf{Examples.} \emph{Same-Side:} a covenant against non-Christian buyers. \emph{Opposite-Side:} a homeowner's own trespass complaint. \emph{Fence-Sitter:} a faith-based bequest condition.

\textbf{Critique.} Read literally, \textit{Shelley} erases the public/private line; Wechsler found no neutral principle, and scholars propose narrower grounds such as willing-party coercion or 42 U.S.C. \S~1982.

\textbf{Key Quotation.} ``That Amendment erects no shield against merely private conduct, however discriminatory or wrongful.'' --- Chief Justice Vinson, 334 U.S. at 13.

\end{document}
